\documentclass[11pt]{article}
\usepackage[a4paper,margin=25mm]{geometry}
\usepackage[T1]{fontenc}
\usepackage{amsmath}
\usepackage{booktabs}
\usepackage{graphicx}
\usepackage{longtable}
\usepackage{listings}
\usepackage{xcolor}
\usepackage{hyperref}

\hypersetup{
  colorlinks=true,
  linkcolor=blue,
  urlcolor=blue,
  citecolor=blue
}
\lstset{
  basicstyle=\ttfamily\small,
  breaklines=true,
  columns=fullflexible,
  keepspaces=true,
  showstringspaces=false
}
\title{Debian based Computer Literacy\\Setup and Everyday Reference}
\author{Hugh Rodgers}
\begin{document}
\maketitle

\section{Install VS Code or VSCodium}
Choose one editor. For VS Code, download the Debian package from the
\href{https://code.visualstudio.com/docs/setup/linux}{official Linux setup
guide} and install it with the graphical software centre or the command shown
there or run sudo apt install ./vscode.deb. 
For VSCodium, follow the\href{https://vscodium.com/#install}{official installation options}; choose
the package source for your distribution. On a managed computer, follow
university policy before installing software.

\section{Install the shared tools}
On Debian or Ubuntu, update the package list and install Python, its virtual
environment and package tools, and Git:
\begin{lstlisting}
sudo apt update
sudo apt install git python3 python3-venv python3-pip
\end{lstlisting}
For a local LaTeX installation, add the compiler and workshop packages:
\begin{lstlisting}
sudo apt install latexmk texlive-latex-recommended texlive-latex-extra texlive-pictures
\end{lstlisting}
Review package-manager changes before confirming. Do not make system changes
on a managed university computer without permission. These package names and
commands are for Debian-based systems; other distributions use different
package managers. See the
\href{https://www.debian.org/doc/manuals/debian-reference/ch02.en.html}{Debian
package-management documentation},
\href{https://www.python.org/downloads/}{Python downloads},
\href{https://git-scm.com/download/linux}{Git for Linux}, and
\href{https://tug.org/texlive/}{TeX Live documentation}.
Overleaf is a browser-based alternative to installing TeX locally.

\begin{lstlisting}
python3 --version
git --version
pwd
\end{lstlisting}

\tableofcontents
\newpage

\section{Terminal, projects, and paths}
The terminal runs a shell, which accepts commands to navigate folders and
start programs. A project is an organised directory; relative paths begin
from the current working directory. Linux filenames are case-sensitive, and
hidden names begin with a dot, such as \texttt{.venv} and \texttt{.git}.
Quote paths that contain spaces.

\begin{lstlisting}
pendulum-lab/
  README.md
  requirements.txt
  data/
    raw_measurements.csv
  analysis/
    analysis.py
    analysis.ipynb
  figures/
  report/
    report.tex
\end{lstlisting}

\begin{longtable}{p{0.29\linewidth}p{0.64\linewidth}}
\toprule Task & Bash command\\ \midrule
Current directory & \texttt{pwd}\\
List files, including hidden & \texttt{ls -la}\\
Move to the parent directory & \texttt{cd ..}\\
Enter a path with spaces & \texttt{cd "My Project"}\\
Copy / move a file & \texttt{cp source.csv copy.csv} / \texttt{mv copy.csv data.csv}\\
Find CSV files & \texttt{find . -name '*.csv'}\\
Search file text & \texttt{grep -n 'time' notes.txt}\\
Open the current folder & \texttt{xdg-open .}, on a desktop system\\
Open the folder in the editor & \texttt{code .} or \texttt{codium .}\\
\bottomrule
\end{longtable}
Inspect paths before copying or moving files: an existing destination may be
overwritten. Prefer the file manager's Trash when removing files.
Command behaviour is documented in the
\href{https://www.gnu.org/software/coreutils/manual/coreutils.html}{GNU
Coreutils manual} and the
\href{https://www.gnu.org/software/findutils/manual/html_mono/find.html}{GNU
Findutils manual}.

\section{Python and project environments}
Create one virtual environment in the project directory containing
\texttt{requirements.txt}. Install project packages inside it, not into the
operating system's Python:
\begin{lstlisting}
python3 -m venv .venv
source .venv/bin/activate
python -m pip install -r requirements.txt
python -m pip check
python -c "import sys; print(sys.executable)"
deactivate
\end{lstlisting}
The printed interpreter path should include \texttt{.venv}. A virtual
environment keeps project packages separate; it is not a security sandbox,
and it should not be committed or copied to another computer. See the
\href{https://docs.python.org/3/library/venv.html}{Python virtual-environment
documentation}. If creating one reports a missing \texttt{ensurepip}, install
the matching distribution package with administrator permission. Do not use
\texttt{sudo pip} or bypass the distribution's externally-managed Python
protection.

\section{Git and LaTeX}
Git works locally; GitHub hosts online copies and collaboration. Saving a file
is not the same as committing it, and a commit is not the same as pushing it.
For the shared commands and safe workflow, see the
\href{../../Git/Git.pdf}{Git and GitHub guide}.

LaTeX source (\texttt{.tex}) is compiled into a PDF. Use Overleaf in a browser
or compile locally with the TeX Live installation above. From the repository
root, build this handout with:
\begin{lstlisting}
latexmk -cd -pdf -interaction=nonstopmode -halt-on-error Linux/Debian_Bassed/Linux.TEX
\end{lstlisting}
The shared
\href{../../LaTeX/LaTeX.pdf}{LaTeX guide} covers report structure, equations,
figures, and references; see also the
\href{https://www.overleaf.com/learn/latex/Learn_LaTeX_in_30_minutes}{Overleaf
LaTeX introduction} and
\href{https://www.latex-project.org/help/documentation/}{LaTeX Project
documentation}.

\section{Shortcuts and Bash configuration}
\begin{itemize}
\item \texttt{Ctrl+C} interrupts the running command; \texttt{Ctrl+D} ends
input or exits an empty shell.
\item \texttt{Tab} completes paths; \texttt{Up} recalls command history;
\texttt{Ctrl+R} searches Bash history.
\item \texttt{Ctrl+Shift+C/V} copies and pastes in many Linux terminal
applications. Desktop shortcuts can vary.
\end{itemize}

Bash reads \texttt{\textasciitilde/.bashrc} for interactive non-login
terminals. To add personal shortcuts, open that file in a text editor and
append simple aliases such as:
\begin{lstlisting}
alias ll='ls -lah'
alias croot='cd ~'
\end{lstlisting}
Save the file, then load the changes in the current terminal:
\begin{lstlisting}
source ~/.bashrc
\end{lstlisting}
Back up and review your own file before editing it. Never replace it wholesale
with another person's configuration. See the
\href{https://www.gnu.org/software/bash/manual/html_node/Bash-Startup-Files.html}{Bash
startup-file documentation}.

\section{Troubleshooting}
\begin{description}
\item[Command not found:] Check that the tool is installed and reopen the
terminal after installation.
\item[Permission required:] Check university policy; do not use administrator
privileges as a general fix.
\item[Python package installation is blocked:] Create and activate the
project's virtual environment.
\item[File not found:] Check \texttt{pwd}, filename case, and the relative path.
\item[No graphical window over SSH:] The remote machine may not have a display.
\item[LaTeX package missing:] Use the TeX distribution's package manager or
Overleaf; an editor extension is not a TeX distribution.
\end{description}

\noindent For Python, fitting, and data-analysis exercises, use the
\href{../../Example_Code/Exersize.ipynb}{shared exercise notebook}.
\end{document}
